% from ega2-1-fr.tex
\oldpage[II]{19}
montrer que $\mathcal{E}$ s'identifie au sous-Module de
$\mathbf{S}(\mathcal{E})$ tel que, \emph{pour tout ouvert $U\subset S$,
toute section $z$ de ce sous-$\mathcal{O}_S$-Module au-dessus de $U$ et tout
$S$-préschéma $X$, l'application $h\mapsto h(z)$ de
$\operatorname{Hom}_{\mathcal{O}_S|U}(\mathbf{S}(\mathcal{E})|U,
\mathcal{A}(X)|U)$ dans $\Gamma(U,\mathcal{A}(X))$ soit un homomorphisme de
$\Gamma(U,\mathcal{A}(X))$-modules}.

Il est immédiat en effet que $\mathcal{E}$ possède cette propriété~; pour
démontrer la réciproque, on peut se borner à prouver que lorsque
$S=\operatorname{Spec}(A)$, $\mathcal{E}=\widetilde{M}$, une section $z$ de
$\mathbf{S}(\mathcal{E})$ au-dessus de $S$ qui (pour $U=S$) possède la
propriété énoncée ci-dessus, est nécessairement une section de
$\mathcal{E}$~; on a alors $z=\sum_{n=0}^{\infty}z_n$, où
$z_n\in\mathbf{S}_n(M)$, et il s'agit de prouver que $z_n=0$ pour $n\ne1$.
Posons $B=\mathbf{S}(M)$ et prenons pour $X$ le préschéma
$\operatorname{Spec}(B[T])$, où $T$ est une indéterminée. L'ensemble
$\operatorname{Hom}_{\mathcal{O}_S}(\mathbf{S}(\mathcal{E}),
\mathcal{A}(X))$ s'identifie à l'ensemble des homomorphismes d'anneaux
$h:B\to B[T]$ (\hyperref[I.1.3.13-fr]{I, 1.3.13}), et d'après ce qu'on a vu
plus haut, on a
$(T.h)(z)=\sum_{n=0}^{\infty}T^nh(z_n)$~: l'hypothèse sur $z$ implique que
l'on a
$\sum_{n=0}^{\infty}T^nh(z_n)=T.\sum_{n=0}^{\infty}h(z_n)$ pour tout
homomorphisme $h$. Si on prend en particulier pour $h$ l'injection
canonique, il vient
$\sum_{n=0}^{\infty}T^nz_n=T.\sum_{n=0}^{\infty}z_n$, ce qui entraîne bien
la conclusion $z_n=0$ pour $n\ne1$.
\end{env}

\begin{proposition}[1.7.15]
\label{II.1.7.15-fr}
Soit $Y$ un préschéma dont l'espace sous-jacent est noethérien, ou un schéma
quasi-compact. Tout $Y$-schéma affine $X$ de type fini sur $Y$ est
$Y$-isomorphe à un sous-$Y$-schéma fermé d'un $Y$-schéma de la forme
$\mathbf{V}(\mathcal{E})$, où $\mathcal{E}$ est un $\mathcal{O}_Y$-Module
quasi-cohérent de type fini.
\end{proposition}

\begin{proof}
En effet, la $\mathcal{O}_Y$-Algèbre quasi-cohérente $\mathcal{A}(X)$ est de
type fini (\hyperref[II.1.3.7-fr]{1.3.7}). Les hypothèses faites entraînent
que $\mathcal{A}(X)$ est engendrée par un sous-$\mathcal{O}_Y$-Module
quasi-cohérent de type fini $\mathcal{E}$
(\hyperref[I.9.6.5-fr]{I, 9.6.5})~; par définition, cela entraîne que
l'homomorphisme canonique
$\mathbf{S}(\mathcal{E})\to\mathcal{A}(X)$ prolongeant canoniquement
l'injection $\mathcal{E}\to\mathcal{A}(X)$ est \emph{surjectif}~; la
conclusion résulte alors de (\hyperref[II.1.4.10-fr]{1.4.10}).
\end{proof}

\section{Spectres premiers homogènes}
\label{section:II.2-fr}

\subsection{Généralités sur les anneaux et modules gradués.}
\label{subsection:II.2.1-fr}

\begin{notation}[2.1.1]
\label{II.2.1.1-fr}
Étant donné un anneau $S$ \emph{gradué} en degrés positifs, nous désignerons
par $S_n$ la partie de $S$ formée des éléments homogènes de degré $n$
($n\geq0$), par $S_+$ la somme (directe) des $S_n$ tels que $n>0$~; on a
$1\in S_0$, $S_0$ est un sous-anneau de $S$, $S_+$ un idéal gradué de $S$ et
$S$ est somme directe de $S_0$ et de $S_+$. Si $M$ est un \emph{module
gradué} sur $S$ (à degrés positifs ou négatifs), on désignera de même par
$M_n$ le $S_0$-module formé des éléments homogènes de degré $n$ de $M$ (avec
$n\in\mathbf{Z}$).

Pour tout entier $d>0$, on désigne par $S^{(d)}$ la somme directe des
$S_{nd}$~; en considérant les éléments de $S_{nd}$ comme homogènes de degré
$n$, les $S_{nd}$ définissent sur $S^{(d)}$ une structure d'anneau gradué.

Pour tout entier $k$ tel que $0\leq k\leq d-1$, on désignera par
$M^{(d,k)}$ la somme directe
